\part{集合、简易逻辑与不等式初步}
\chapter{数理逻辑}
\section{集合}
\subsection{集合的定义、特性与表示}
在数学学习中，我们会研究一些具有共同特征的对象，我们往往称之为元素.
比如所有的质数，过某个定点的直线束......此时我们会把这些元素放在一起，得到一个集合.
集合及相关理可认为是现代数学的根基，是拥有一套公理的严格数学语言.
但是我们在高中并不会接触晦涩的公理，而是着重研究一些实用，直观的对象.
\begin{definition}
    集合(set)是一些不同的、能够确定的研究对象汇集在一起构成的数学结构，其中研究的对象称之为元素(element)。
\end{definition}

我们常用大写字母$A,B,C...$表示集合，一些特殊的集合可能会使用特殊字体如$\mathcal{B},\mathbb{R},\mathfrak{F} $表示。
相对地，我们会使用小写字母$a,b,c...$来表示集合中的元素.

若元素$a$在集合$A$中，则记作\[a \in A\]

相对地，若元素$a$不在集合$A$中，则记作\[a \notin A\]

当然，如果有这样的一个集合，没有任何元素属于这个集合，那么我们称之为空集(empty set/null set)，符号为$\varnothing $.

对于两个集合$A,B$，若构成两个集合的元素相同，则记作$A,B$两个集合相等，即\[A=B\]
\begin{example}
    判断下面叙述的对象能否构成集合。若能构成集合，试图列举一些集合中的元素.
    
    $(1)$教室中比较高的同学，教室中最高的同学.

    $(2)$西游记中所有的妖怪.

    $(3)$所有的奇数，所有的偶数，所有的素数.

    $(4)$在平面上，到$A$$B$两点距离相同的所有的点.

    $(5)$单词success中的所有字母.

    $(6)$不能表示为整数或整数之比的有理数.

    $(7)$平面上所有的平行四边形.

    $(8)$除以$2$余$1$的所有整数.
\end{example}

\begin{definition}[子集]
    \[A\subseteq  B\]

    \[A \subsetneqq B\]
\end{definition}
\begin{theorem}
    \[A=B \Leftrightarrow A\subseteq B \text{且} B\subseteq A \]
\end{theorem}



知识提纲
\\集合元素的三个特性：
\\集合与元素的关系，集合与集合的关系
\\子集，空集，真子集，非空真子子集
\\常见的集合：$\mathbb{N},\mathbb{Z},\mathbb{Q},\mathbb{R},\mathbb{C}$等
\\含有n个元素的集合共有多少个子集？多少真子集？多少非空真子集？
\\集合间的运算与交并补，使用Venn图表示

\begin{example}
    \[\left\{x|y=x^2+2x+2,x\in \mathbb{R} \right\}\]
\[\left\{y|y=x^2+2x+2,x\in \mathbb{R} \right\}\]
\[\left\{(x,y)|y=x^2+2x+2,x\in \mathbb{R} \right\}\]
\[\left\{y=x^2+2x+2 \right\}\]
\end{example}

提示：字母究竟是什么意思？

\begin{example}
    \[
    \emptyset \in \left\{\emptyset\right\},
    \emptyset \subseteq \left\{\emptyset\right\},
    \emptyset \notin \emptyset,
    \emptyset \subseteq \emptyset
\]
\end{example}

\begin{example}集合A为$\left\{1,2,3,4,5,6\right\}$，集合为$\left\{4,5,6,7,8\right\}$。
则满足$S\subseteq A$且$S\cap B\neq \emptyset$的S有几个？
    
\end{example}
解：提示：将元素分为123,456,78三类。


特殊的集合表示法：区间
\\用法细则：
\\不等式的解集写集合，形如$\left\{x|-3<x<2,x \in \mathbb{R} \right\}$
\\单调区间写区间
\\取值范围尽量写区间，不要写一个光秃秃的不等式
\begin{example}
    已知 $ a \in \mathbb{R} $, 关于  $x$ 的方程 $ x^{2}+a=x  $的解组成的集合为 $ A(A \neq   \varnothing)$
    $,\left(x^{2}+a\right)^{2}+a=x $ 的解组成的集合为 $ B .$

$(1) $求证  $: A \subseteq B ;$

$(2)$ 若$  A=B$ , 求实数 $ a $ 的取值范围.
\end{example}


\clearpage
\subsection{集合的运算}
\clearpage
\hw
\begin{homework}
    \ 

    定义$1$：称一个以集合作为元素的集合为族$(collection)$.

    定义$2$：称非空集合$X$的一个子集族$\Gamma$为X上的一个拓扑$(topology)$是指：\\
\ding{172}$ \emptyset,X\in \Gamma$;\\
\ding{173}$\Gamma$的任意多元素的并集仍在$\Gamma$中;\\
\ding{174}$\Gamma$的任意有限多元素的并集仍在$\Gamma$中.

    $(1)$族$P=\{\emptyset,X\}$，族$Q=\{A|A\subseteq X\}$,判断$P$、$Q$是否是$X$上的一个拓扑；
    
    $(2)$$\Gamma$为有限集$X$上的一个拓扑，证明：由\(\Gamma\)的所有元素关于$X$的补集构成的族$\Gamma^c$也为$X$上的一个拓扑;

    $(3)$证明：$X$为无限集合，$\tau_f=\{\complement_X A|A\text{为}X\text{的有限子集}\}\bigcup\{\varnothing\}$，则$\tau_f$为$X$上的一个拓扑.
\end{homework}



\clearpage

\section{简易逻辑}
\subsection{从集合到命题}
命题最最基础的定义是“可以判断真假的陈述句。”，在一定的范围内可以判定真假是判断是否是命题的核心方法。
一个误区是所有的命题都可以写成若p则q的形式，这是不对的。实际上命题的核心形式是$x \in \mathbb{X} $,很显然元素属于集合与否是容易判断的。
但是我们见到的大多数命题都经过了各种修饰与复合，使得看起来和集合已经没有什么关系了。
我们最常见的两种命题就是若p则q和带有$\forall,\exists $的命题。（当然两者结合在一起也可以）
\clearpage
\subsection{四种命题*}

\clearpage
\subsection{充分条件，必要条件，充要条件}
\begin{definition}
    若“若p，则q”是真命题，则记作$p \Rightarrow q$，称p是q的充分条件，q是p的必要条件。
反之，若“若p，则q”是假命题，则记作$p \not\Rightarrow  q$，称p不是q的充分条件，q不是p的必要条件。
\end{definition}

思考：充分一词很容易理解，那么必要一词如何理解？即q不成立，p一定不成立。请用一些具体的命题来验证一下这个说法。
\begin{definition}
    若满足$p \Rightarrow q$且$q \Rightarrow p$，则称p与q互为充分必要条件，简称充要条件。
记作
\[p \Leftrightarrow q\]
\end{definition}

\begin{exercise}
    请写出“p是q的充分不必要条件”和“p是q的必要不充分条件”的等价命题。
\end{exercise}


与集合的关系
\\若$p \Rightarrow q$，其中$p$是$x \in \mathbb{P} $，$q$是$x \in \mathbb{Q} $,
则$\mathbb{P} $与$\mathbb{Q} $的关系是？（$\mathbb{P}\subseteq\mathbb{Q} $）简记为：在小范围里一定在大范围里

例8
\begin{example}
    已知集合$\mathbb{P} =\{x|a-4<x<a+4\}$，$\mathbb{Q} =\{x|1<x<3\}$,“$x \in \mathbb{P} $”是“$x \in \mathbb{Q} $”
的必要条件，则实数a的取值范围是？
\end{example}

\begin{jie}
    “$x \in \mathbb{P} $”是“$x \in \mathbb{Q} $”的必要条件，即“$x \in \mathbb{Q} $”是“$x \in \mathbb{P} $”
的充分条件，根据上文可得$\mathbb{Q} \subseteq \mathbb{P} $.列出不等式组

\[  
    \left\{
    \begin{lgathered}
    a-4\leq 1\\
    3\leq a+4\\
    \end{lgathered}
    \right.
\]
最后解得$a\in [-1,5]$
\end{jie}
\begin{exercise}
    判断一下说法的正误：
    \begin{enumerate}
        \item $a=b$是$ac=bc$的充分不必要条件；
        \item $\frac{1}{a}>\frac{1}{b}$是$a<b$的既不充分也不必要条件；
        \item $a\neq 0$是$ab \neq 0$的必要不充分条件；
        \item $a>b>0$是$a^n>b^n(n\in \mathbb{N}\text{且}n\leq2)$的充要条件。
    \end{enumerate}
\end{exercise}
\begin{exercise}
    已知条件p:$|x-1|\leq 2$，条件q:$x>a$，且满足q是p的必要不充分条件，则a取值为？
\end{exercise}

\subsection{全称量词、存在量词与命题的否定}
\begin{definition}
    在逻辑中，诸如“所有的”，“任意一个”的短语叫做全称量词，用记号$\forall$表示；
同理，诸如“存在”“至少一个”的短语叫做存在量词（特称量词），用记号$\exists$表示。记作
\[\forall x\in M,P(x).\exists x\in M,P(x).\]
其含义也可以看做是
\[\forall x\in M,x\in \{x|P(x)\}.\exists x\in M,x\in \{x|P(x)\}.\]
\end{definition}
翻译条件
\begin{example}
    $\forall x\in A,\exists y\in B : y=x$的等价条件是？
\end{example}

\begin{jie}
    \[A\subseteq B\]
\end{jie}


“你有多大脚，我有多大鞋。”


练习：借助函数图像，使用最大值与最小值的语言翻译下列叙述。
\\$\forall x_1 \in D_1,\forall x_2 \in D_2,f(x_1)>g(x_2).\Leftrightarrow$
\\$\exists x_1 \in D_1,\exists x_2 \in D_2,f(x_1)>g(x_2).\Leftrightarrow$
\\$\forall x_1 \in D_1,\exists x_2 \in D_2,f(x_1)>g(x_2).\Leftrightarrow$
\\$\exists x_1 \in D_1,\forall x_2 \in D_2,f(x_1)>g(x_2).\Leftrightarrow$
\\$\forall x \in D$,使得$f(x)>g(x)$恒成立$\Leftrightarrow$
\\$\exists x \in D$,使得$f(x)>g(x)$可以成立$\Leftrightarrow$
\\$\forall x_1 \in D_1,\exists x_2 \in D_2,f(x_1)=g(x_2).\Leftrightarrow$
\\$\exists x_1 \in D_1,\exists x_2 \in D_2,f(x_1)=g(x_2).\Leftrightarrow$


首先应当引入一个概念——排中律，即一个命题与它的否定应该有且仅有一个成立，
有时候我们也成命题的否定为原命题的“非命题”，这写都意味着原命题与命题的否定必然是一对一错的关系。
而我们主要研究的就是带有量词的。
\begin{theorem}
    随着$\neg$划过整个命题，命题中的“$\exists$”和$“\forall$”互换，最终停留在表达式前面。
\end{theorem}

\begin{example}
    \[\forall x\in M,P(x).\exists x\in M,P(x).\]
对应的否定是
\[\exists x\in M,\neg P(x).\forall x\in M,\neg P(x).\]
\end{example}


多个量词自然也不难，应用上一节的练习验证这个结论。

思考：全程量词和存在量词的否定的结果是彼此互换的逻辑是什么？是否可以帮助我们证明或否定一些含有量词的命题？

前面量词部分的转换是容易的，重点还是后面从$P(X)$到$\neg P(x)$的做法。
\\在前文我们提到过，逻辑用语的最本质形式是$x \in \mathbb{X} $，
而我们要做的就是把$\in$改为$\notin$即可。相应的，也就是将$=$改为$\neq$,$>$改为$\leq$

思考：若p，则q的形式的命题如何改成否定形式呢？

\clearpage
\chapter{不等式初步}
\section{不等式}
\subsection{不等式的性质}
研究不等式，其实就是在追溯一个古老的问题———比大小.我们把目光主要放在$\mathbb{R}$上，由于这个集合良好的性质
其中的任意两个实数$a,b$，一定是可以比较大小的，也就是它们一定满足
\[a>b,a=b,a<b\]
其中之一。注意，不是所有数集中的元素都可以比较大小的，一个经典的例子就是复数集$\mathbb{C}$，这些会在
后面学习中学到.

话说回来，我们如何定义$a$与$b$两个数（代数式）的大小呢？
我们已经有了“正数>$0$>负数”的这样的结论，所以两数的大小只要归结到两数之差和$0$的即可，
这是基于我们过去学习的一个直观的想法.
\begin{definition}
    \ 

    $a>b$是指$a-b>0$，即$a>b\Leftrightarrow a-b>0$.

    同理可以定义$<,\geq,\leq,\neq$.

\end{definition}
由此定义我们可以推导出来不等式的若干性质，或许在初中数学中已经了解过它们了，但是我们这里要看看如何严格地利用定义证明这些性质。
\begin{character}
    \ 

    $(1)$反对称性\[a>b\Leftrightarrow b<a\]

    $(2)$传递性\[a>b,b>c \Rightarrow a>c\]

    $(3)$可加性\uppercase\expandafter{\romannumeral1} \[\forall c,a>b\Leftrightarrow a+c>b+c\]

    $(4)$可加性\uppercase\expandafter{\romannumeral2}\[a>c,b>d \Rightarrow a+b>c+d \]

    $(5)$可乘性\uppercase\expandafter{\romannumeral1}
\[\forall c>0,a>b\Leftrightarrow ac>bc\]
\[\forall c<0,a>b\Leftrightarrow ac<bc\]

    $(6)$可乘性\uppercase\expandafter{\romannumeral2}
\[a>b>0,c>d>0 \Rightarrow ac>bd\]

    $(7)$可乘方性
\[a>b>0 \Rightarrow a^n>b^n (n \in \mathbb{N} ,n\geq 2)\]

    $(8)$可开方性
\[a>b>0 \Rightarrow \sqrt[n]{a}>\sqrt[n]{b} (n \in \mathbb{N} ,n\geq 2)\]

\end{character}
\begin{proof}
    \ 

    仅证明$(4)(6)(7)(8)$，其余作为练习.
\end{proof}

练习 \[a>b\Rightarrow \frac{1}{a}\_\_\frac{1}{b}\]
例1 一个很抽象高考不会考但是得会的题目
\\已知$-3\leq a+b\leq2,-1\leq a-b\leq 4$，求$a,b,3a-2b$的取值范围。
\\解：求$a,b$的取值范围无疑是很简单的，但是在求$3a-2b$的取值范围是时候是否可以直接套用呢？答案是不能
答案是不可以。
\\两式相加得$-4\leq 2a \leq 6$，即$-2\leq a \leq 3$.同理两式相减（注意不等式的减法的算法）
得到$-7\leq 2b\leq 3$，即$-3.5\leq b\leq 1.5$.
\\在求$3a-2b$的时候，使用最初的不等式配凑的带$3a-2b$,而不要用a,b的取值范围直接相加减得到。
使用直接相加减会得到$-9\leq 3a-2b\leq16$,这个16是哪里来的？应该是a取最大值3，b取最小值-3.5的时候得到的。
但是此时$a-b=7.5$,并不符合题目最初的式子的要求，说明随着多次使用同向不等式的加减法，逐渐扩大了最初的取值范围。
因此才要使用所谓的配凑法。
令\[3a-2b=m(a+b)+n(a-b)=a(m+n)+b(m-n)\]
对比等号两端，得到方程组
\[  
    \left\{
    \begin{lgathered}
    m+n=3\\
    m-n=-2
    \end{lgathered}
    \right.
\]
解得
\[  
    \left\{
    \begin{lgathered}
    m=0.5\\
    n=2.5
    \end{lgathered}
    \right.
\]
即\[3a-2b=0.5(a+b)+2.5(a-b)\],乘上对应倍数相加后得到
\[-4\leq 3a-2b\leq11\]
很明显这个范围比之前缩小了不少。等学完直线方程之后再来看这个就会清晰很多了。这个东西本质上是新高考中被删掉的线性规划。
\subsection{比大小与证明不等式}
比大小和不等式的研究方法有许多，我们手里的工具现在还很少，我们会先介绍一些简单的方法和例子.
\begin{theorem}\label{bidaxiao}
    \[\forall a,b\in\mathbb{R},a>b\Leftrightarrow a-b>0\]
    \[\forall a,b\in\mathbb{R}_+,a>b\Leftrightarrow \frac{a}{b}>1\]
\end{theorem}

由不等式的性质，定理$\ref*{bidaxiao}$是显然的。也就是说对于绝大多数的比大小问题，
归结于作差或者是作商，再判断差或商的性质。针对不同的结构，我们需要灵活选择方法.
\begin{example}
    \ 

    $(1)$比较$a^2+b^2$与$2(2a-b)-5$

    $(2)$$\forall a,b\in\mathbb{R}_+$,比较
    \[\frac{a}{b^2}+\frac{b}{a^2}\text{与}\frac{1}{a}+\frac{1}{b}\]
\end{example}
\begin{jie}
    \ 

    $(1)$作差得
    \begin{align*}
    &\quad\  a^2+b^2-(2(2a-b)-5)\\
    &=a^2+b^2-4a+2b+5\\
    &=a^2-4a+4+b^2+2b+1\\
    &=(a-2)^2+(b+1)^2 \geq 0
    \end{align*}
    
    故$a^2+b^2\geq 2(2a-b)-5$，当且仅当$a=2$且$b=-1$时取得等号.

    $(2)$
    \begin{way1}
    作差得
    \[\frac{a}{b^2}+\frac{b}{a^2}-(\frac{1}{a}+\frac{1}{b})
    =\frac{a^3+b^3-ab^2-a^2b}{a^2b^2}=\frac{(a-b)^2(a+b)}{a^2b^2}\geq0\]

    第二个等号处的因式分解只要将四项两两重新组合即可。这里$\forall a,b\in\mathbb{R}_+$保证了
    最后的式子每一个因式都是非负的，故整个式子是非负的.
    \end{way1}
    \begin{way2}
    显然两个式子都为正，故做商得
    \begin{align*}
    &\quad\ \frac{\dfrac{a}{b^2}+\dfrac{b}{a^2}}{\dfrac{1}{a}+\dfrac{1}{b}}\\
    &=\frac{a^3+b^3}{ab^2+a^2b}\\
    &=\frac{(a+b)(a^2-ab+b^2)}{ab(a+b)}\\
    &=\frac{a^2-ab+b^2}{ab}\\
    &=\frac{a}{b}+\frac{b}{a}-1 \geq 2\sqrt{\frac{a}{b}\cdot \frac{b}{a}}-1=1
    \end{align*}

    当且仅当满足$\dfrac{a}{b}=\dfrac{b}{a}$即$a=b$时取等.
    \end{way2}
\end{jie}
\begin{example}
    \ 

    $(1)$已知$x>0$，求证\[\sqrt{1+x}<1+\dfrac{x}{2}\]
    
    $(2)$已知$x\in[0,1]$，求证\[x^3+\frac{1}{x+1}\geq 1-x+x^2\]

    $(3)$已知$n>0$，比较\[\sqrt{n}-\sqrt{n-1}\text{与}\sqrt{n-1}-\sqrt{n-2}\]
\end{example}





\begin{proposition}
    hh
    
\end{proposition}












\clearpage
\section{常见的不等式}
\subsection{基本不等式}
从完全平方公式$(a-b)^2=a^2-2ab+b^2$出发
得到不等式
\[(a-b)^2=a^2-2ab+b^2 \geq 0 \Leftrightarrow a^2+b^2\geq 2ab.\text{当且仅当$a=b$取等} \]
特别地，$a>0,b>0$时，用$\sqrt{a},\sqrt{b}$替换$a,b$,得到不等式
\[\frac{a+b}{2}\geq \sqrt{ab},\text{当且仅当$a=b$取等}\]

当然，我们也经常使用两侧平方后的形式，即
\[ab\leq \frac{(a+b)^2}{4}\]
对于$\dfrac{a+b}{2}$
\begin{example}
    已知$x,y>0$，求证：

    $(1)$若$xy=P$（定值），当$x=y$时，和$x+y$有最小值$2\sqrt{P}$

    $(2)$若$x+y=P$（定值），当$x=y$时，积$xy$有大值$\frac{1}{4}S^2$
\end{example}
\begin{example}
   \[ \text{当$x>0$时，求}x+\frac{1}{x}\text{的最小值}\]
\end{example}
\begin{jie}
    $x>0$,则
\[x+\frac{1}{x} \geq 2\sqrt{x\cdot\frac{1}{x}}=2\]
\end{jie}

归纳均值不等式的使用步骤：一正，二定，三相等，顺序不能乱
\begin{example}
    \ 

    $(1)$\[x>1,\text{求}x+\frac{1}{x-1}\text{的最小值;}\]

    $(2)$\[ab=8\text{，求$a+2b$的最小值;}\]

    $(3)$\[x>2,\text{求}\frac{x^2-2x+4}{x-2}\text{的最小值}\]

    $(4)$
    \[k\leq 1\text{，求}\frac{x^2+k+1}{\sqrt{x^2+k}}\text{的最小值}\]

    $(5)$\[x,y\text{同号，求}\frac{x}{y}+\frac{y}{x}\text{的最小值}\]

    $(6)$\[\text{求}\frac{xy}{x^2+3xy+y^2}\text{的最大值}\]

    $(7)$\[a,b\in \mathbb{R}_+,\text{求}\frac{a}{a+2b}+\frac{b}{a+b}\text{的最小值}\]
\end{example}
\begin{exercise}
    \ 

    $(1)$\[x\geq1,\text{求}x+\frac{1}{x}+\frac{16x}{x^2+1}\text{的最小值}\]
\end{exercise}
\begin{example}
    \[x\in(0,1),\text{求}x(1-x)\text{的最小值}\]
\end{example}
\begin{jie}
    显然把这题作为二次函数求最值来做是没有任何问题的，但是它也有天然的“和是定值”的优秀形式。
\end{jie}
\begin{example}
    \ 

    $(1)$\[x\in(0,1),\text{求}x(2-3x)\text{的最小值}\]

    $(2)$\[3a^2+2b^2=5,\text{求}(2a^2+1)(b^2+2)\text{的最大值}\]

    $(3)$\[5x^2y^2+y^4=1,\text{求}x^2+y^2\text{的最小值}\]
\end{example}

练习：(1)$x,y$为不相等的正数，求证
\[\frac{2xy}{x+y}<\sqrt{xy}\]

\begin{example}
    \[x>0,y>0,x+y=1,\text{求}\frac{1}{x}+\frac{4}{y}\text{的最小值}\]
\end{example}
\begin{jie}
    
\end{jie}
\begin{example}
    \ 

    $(1)a,b>0,2a+b=1,$
    \[\text{求}\frac{1}{a+1}+\frac{4}{b+1}\text{的最小值}\]

    $(2)a,b>0,a+b=1,$
    \[\text{求}\frac{1}{3a+2}+\frac{1}{3b+2}\text{的最小值}\]

    $(3)x,y>0,\dfrac{2}{x}+\dfrac{1}{y}=2,$
    \[\text{求}x+2y\text{的最小值}\]

    $(4)x,y>0,a+\dfrac{1}{b}=2$
    \[\text{求}\frac{4}{a}+b\text{的最小值}\]
\end{example}
\begin{exercise}
    \ 

    $(1)x,y>0,x+3y=5xy,$\[\text{求}3x+4y+xy\text{的最小值}\]
\end{exercise}
\begin{example}
    \ 

    $(1)a,b>0,ab=a+b+3,$
    \[\text{求}ab\text{的取值范围}\]

    $(2)x,y>0,x+y+xy=3,$
    \[\text{求}x+y\text{的取值范围}\]
\end{example}
\begin{theorem}[n元均值不等式]
    \ 

    $a_1,a_2,\dots,a_n\geq0$,则
    \[\frac{a_1+\dots+a_n}{n}\geq\sqrt[n]{a_1 \dots a_n}\]
    
    当且仅当$a_1=a_2=\dots=a_n$时等号成立。
\end{theorem}
\begin{proof}
    使用向前——向后数学归纳法。
    
    $n=1,2$时,显然成立.下证明若$n=2^k$时不等式成立,则$n=2^{n+1}$时不等式仍然成立.

    假设$n=2^k$时，有\[\]

    接下来证明若$n=k$时不等式成立,则$n=k-1$时不等式仍然成立.
\end{proof}
\begin{example}
    \ 

    求证
    \[\forall a,b,c\in \mathbb{R}\qquad a^2+b^2+c^2\geq ab+bc+ac\]
\end{example}
\begin{proof}
    由\[a^2+b^2\geq 2ab\]
    
    类推可得\[b^2+c^2\geq 2bc\]\[a^2+c^2\geq 2ac\]

    累加得\[2a^2+2b^2+2c^2\geq 2ab+2bc+2ac\]

    约去两侧的$2$即得证。
\end{proof}
\begin{zhu}
    请用分析法写此题的证明过程。
\end{zhu}
\clearpage
\begin{exercise}
    \ 

    $(1)a,b,c,d>0$,求证:
    \[(a+b)(b+c)(c+d)(d+a)\geq 16abcd\]

    $(2)x,y,z>0$,求证：
    \[(\frac{y}{x}+\frac{z}{x})(\frac{x}{y}+\frac{z}{y})(\frac{x}{y}+\frac{y}{z})\geq8\]

\end{exercise}
\clearpage
\subsection{不等式链}
\[\sqrt{\frac{a^2+b^2}{2}} \geq \frac{a+b}{2}\geq \sqrt{ab}\geq \frac{2}{\dfrac{1}{a}+\dfrac{1}{b}}\]

\clearpage
\subsection{糖水不等式}
\[a,b\in \mathbb{R}_+,a>b,\frac{b+m}{a+m}>\frac{b}{a}\]
\[a,b\in \mathbb{R}_+,a<b,\frac{b+m}{a+m}<\frac{b}{a}\]
\[a,b,m\in \mathbb{R}_+,a>b,a>m,\frac{b-m}{a-m}<\frac{a}{b}\]
\[a,b,c,d\in \mathbb{R} ,\frac{b}{a}<\frac{d}{c}\text{则}\frac{b+d}{a+c}>\frac{d}{c}>\frac{b}{a}\]

\clearpage
\subsection{绝对值不等式}
\begin{theorem}[绝对值不等式]
  \[||a|-|b||\leq|a\pm b|\leq |a|+|b|\]
  \[|\sum_{i=1}^{n}a_i|\leq\sum_{i=1}^{n}|a_i|\]  
\end{theorem}
\begin{corollary}
    饿啊
\end{corollary}
\begin{example}
    \ 

    $(1)$$|x-m|+|x-n|$

    $(2)$$|x-m|-|x-n|$

    $(3)$$a|x-m|+b|x-n|$
\end{example}
\begin{example}
    \[f(x)=\sum_{i=1}^{n}|x-a_i|=|x-a_1|+|x-a_2|+\cdots+|x-a_i|\]

    $(1)$当$n=2k-1(k\in \mathbb{N} *)$,在$f(a_k)$处取得最小值.

    $(2)$当$n=2k(k\in \mathbb{N} *)$,在$f(a_k)$到$f(a_{k+1})$取得最小值.
\end{example}
\begin{example}
    \[\text{已知}x,y\text{满足}|x+y|+|x|+|y|=2\]

    求$x,y,x+y,x-y$的取值范围。
\end{example}
\begin{definition}[距离空间]
    
\end{definition}
\begin{axiom}
    vv
\end{axiom}
\clearpage
\subsection{柯西不等式*}
\begin{theorem}[Cauthy不等式]
    \[(a^2+b^2)(c^2+d^2)\geq(ac+bd)^2\]
\[(\sum_{i=1}^{n}a_i^2)(\sum_{i=1}^{n}b_i^2)\geq(\sum_{i=1}^{n}a_ib_i)^2\]
\end{theorem}

\begin{proof}
    \ 

    \begin{way1}
    \[(a^2+b^2)(c^2+d^2)=(ac+bd)^2+(bc-ad)^2\geq (ac+bd)^2\]
    \[(\sum_{i=1}^{n}a_i^2)(\sum_{i=1}^{n}b_i^2)-(\sum_{i=1}^{n}a_ib_i)^2=
    \sum_{1\leq i<j\leq n}^{n}(a_ib_j-a_jb_i)^2\geq 0\]
    \end{way1}
    \begin{way2}
    记\[f(x)=\sum_{i=1}^{n}(a_i+xb_i)^2
    =(\sum_{i=1}^{n}b_i^2)x^2+2(\sum_{i=1}^{n}a_ib_i)x+\sum_{i=1}^{n}a_i^2\geq 0\]

    这等价于二次函数的判别式
    \[\Delta=4(\sum_{i=1}^{n}a_ib_i)^2-4(\sum_{i=1}^{n}a_i^2)(\sum_{i=1}^{n}b_i^2)\leq 0\]
    
    去掉系数$4$即得证明。
    \end{way2}
\end{proof}
简单记忆为“和的积$\geq$ 积的和”，在学完向量相关内容后，会有对柯西不等式更深入的了解。
\begin{corollary}
    \ 

    $(1)$\[\sqrt{(a^2+b^2)}\sqrt{(c^2+d^2)}\geq |ac+bd|\]

    $(2)$\[(a+b)(c+d)\geq(\sqrt{ac}+\sqrt{bd})^2\]

    $(3)$\[\sqrt{(a^2+b^2)}\sqrt{(c^2+d^2)}\geq |ac|+|bd|\]
\end{corollary}
\begin{example}
    \ 

    $(1)$
    \[x+3y=4,\text{求}4x^2+y^2\text{的最小值}\]   

    $(2)$
    \[\text{求}5\sqrt{x-1}+\sqrt{10-2x}\text{的最大值}\]
\end{example}
\begin{jie}
    \ 

    $(1)$
    \[(4x^2+y^2)(\frac{1}{4}+9)\geq(x+3y)^2\]

    $(2)$
    \begin{align*}
    &\quad\ (5\sqrt{x-1}+\sqrt{10-2x})^2 \\
    &=(5\sqrt{x-1}+\sqrt{2}\cdot\sqrt{5-x})^2\\
    &\leq(5^2+(\sqrt{2})^2)(x-1+5-x) \\
    &=108
    \end{align*}
  
\end{jie}
\begin{theorem}[权方和不等式]
    \[a,b,x,y>0,\frac{a^2}{x}+\frac{b^2}{y}\geq\frac{(a+b)^2}{x+y}\]
    
    当且仅当$\dfrac{a}{x}=\dfrac{b}{y}$时取等号.

    \[a_i,b_i>0,\sum_{i=1}^{n}\frac{a_i^2}{b_i}\geq\frac{(\sum_{i=1}^{n}a_i)^2}{\sum_{i=1}^{n}b_i}\]

    当且仅当$\dfrac{a_i}{b_i}=k$时取等.
\end{theorem}
\begin{proof}
    由柯西不等式
    \[\sum_{i=1}^{n}\frac{a_i^2}{b_i}\sum_{i=1}^{n}b_i\geq(\sum_{i=1}^{n}a_i)^2\]
\end{proof}
\begin{example}
    求证\[\frac{a^2}{(a+b)(a+c)}+\frac{b^2}{(b+a)(b+c)}+\frac{c^2}{(c+a)(c+b)}\geq \frac{3}{4}\]
\end{example}
\begin{proof}
        \ 
        
        由权方和不等式得
        \begin{align*}
        \text{原式}&\geq \frac{(a+b+c)^2}{(a+b)(a+c)+(b+a)(b+c)+(c+a)(c+b)}\\   
                 &=\frac{a^2+b^2+c^2+2ab+2bc+2ac}{a^2+b^2+c^2+3ab+3bc+3ac}\\
                 &=1-\frac{ab+bc+ac}{a^2+b^2+c^2+3ab+3bc+3ac}
        \end{align*}
        
        第一个不等号的取等条件是
        
        而
        \begin{align*}
        &\quad\ \frac{ab+bc+ac}{a^2+b^2+c^2+3ab+3bc+3ac}\\
        &=\frac{2ab+2bc+2ac}{2a^2+2b^2+2c^2+6ab+6bc+6ac}\\
        &=\frac{2ab+2bc+2ac}{a^2+b^2+b^2+c^2+a^2+c^2+6ab+6bc+6ac}\\
        &\leq \frac{2ab+2bc+2ac}{2ab+2bc+2ac+6ab+6bc+6ac}\\
        &=\frac{1}{4}
        \end{align*}
        
        故原式$\geq \dfrac{3}{4}$，当且仅当$a=b=c$时取等号.
        
\end{proof}
\clearpage
\section{常见不等式的解法}
\subsection{解绝对值不等式}
零点分段法
\begin{example}
    用零点分段法解下列绝对值不等式。

    $(1)\quad|2x+1|\leq 3$

    $(2)\quad|x-1|+|x+1|>4$

    $(3)\quad|x+3|-|2x+1|>-4$
\end{example}
\begin{example}
    求下列不等式
\end{example}
\clearpage
\subsection{解一元二次不等式、分式不等式与特殊的高次不等式}
很正常的解方程：画图
\\带绝对值的；翻折
\\双向不等式：多画一条线
\begin{example}解一元二次不等式
    \[0 \leq x^2-2x-3<5\]
\[|x^2-5x+6|<1\]
\end{example}


步骤：移项通分
\begin{example}解分式不等式
    \[\frac{x-1}{x+2} \leq 0\]
\[\frac{x-1}{2x+2} \leq 2\]
\end{example}


画图，数轴标根，奇穿偶不穿
\begin{example}
    \[(x-1)^2(x-5)^3(x+1)(x+7)^4 \geq 0\]
\end{example}



\clearpage
